\subsection{Queries}

\par
Before running any queries, the four datasets (ingested in Phase 2) had to be
re-loaded and registered as permanent Spark SQL tables. This meant using either
\texttt{saveAsTable} or saving using parquet within \texttt{spark.sql}. I chose
the former.

\par
My first attempt at \texttt{saveAsTable} ran without errors, but the tables
came back empty. Running \texttt{DESCRIBE FORMATTED} showed why: since I
hadn't set an explicit path, Spark defaulted to
\texttt{spark.sql.warehouse.dir}, which pointed at a local filesystem path
(\texttt{file:/opt/spark/spark-warehouse/...}). On a multi-container cluster,
``local'' means local to whatever worker ran the write, so the files ended up
on that worker's own disk, invisible to everyone else. I fixed it by passing
an explicit HDFS path via the \texttt{path} option on each write, so the table
data lives somewhere every container can actually reach:

\begin{figure}[H]
	\centering
	\includegraphics[width=0.9\textwidth]{4_1_hdfs_write_tables.png}
	\caption{Registering the four permanent tables with explicit HDFS paths under
		\texttt{/practice/tables/}}
\end{figure}

\par
I confirmed the tables being registered under the \texttt{practice} database
and being permanent with \texttt{isTemporary = false} with \texttt{SHOW TABLES}:

\begin{figure}[H]
	\centering
	\includegraphics[width=0.9\textwidth]{4_1_show_tables_practice.png}
	\caption{\texttt{SHOW TABLES IN practice} \textemdash\ countries, airlines,
		airports, and routes all registered as permanent tables}
\end{figure}

\par
For each of the following queries, I used both the DataFrame API and an
equivalent SQL statement against the permanent tables, and confirmed the two
approaches returned identical results.

\subsubsection*{Q1: How many flights depart from Brazil?}

\par
Filtering \texttt{routes} on \texttt{source\_country = 'Brazil'} and counting
gives 1304 flights, matching between the DataFrame API and SQL:

\begin{figure}[H]
	\centering
	\includegraphics[width=0.9\textwidth]{4_1_Q1.png}
	\caption{Q1 result: 1304 flights depart from Brazil (DataFrame API and SQL)}
\end{figure}

\subsubsection*{Q2: Which airports are located in the NW quadrant?}

\par
This filters the \texttt{hemisphere\_quadrant} column (derived in Phase 2 via
UDF) for the value \texttt{NW}. Results are airports with positive latitude and
negative longitude, such as those across Greenland, Iceland, and Canada:

\begin{figure}[H]
	\centering
	\includegraphics[width=0.9\textwidth]{4_1_Q2.png}
	\caption{Q2 result: airports in the NW quadrant (DataFrame API and SQL)}
\end{figure}

\subsubsection*{Q3: Which countries have the most airports above 2000 feet
	altitude?}

\par
Filtering \texttt{airports} on \texttt{altitude > 2000}, grouping by country,
and ordering by count descending. The United States leads with 234 airports
above 2000 feet, followed by China (76), South Africa (66), and Brazil (63):

\begin{figure}[H]
	\centering
	\includegraphics[width=0.9\textwidth]{4_1_Q3.png}
	\caption{Q3 result: countries ranked by number of airports above 2000 feet
		altitude}
\end{figure}

\subsubsection*{Q4: What are the top 10 routes with the most stops?}

\par
Ordering \texttt{routes} by \texttt{stops} descending and taking the top 10.
Note that the OpenFlights \texttt{stops} field in this dataset is effectively
binary in practice \textemdash\ the maximum value present is 1, so the
``top 10'' result is a set of 10 one-stop routes rather than a strict ranking
by increasingly large stop counts:

\begin{figure}[H]
	\centering
	\includegraphics[width=0.9\textwidth]{4_1_Q4.png}
	\caption{Q4 result: top 10 routes by number of stops (all tied at 1 stop)}
\end{figure}

\subsubsection*{Q5: What is the average number of routes per airline?}

\par
Grouping \texttt{routes} by \texttt{airline\_id}, counting routes per airline,
then averaging across airlines gives 127.42 routes per airline on average:

\begin{figure}[H]
	\centering
	\includegraphics[width=0.9\textwidth]{4_1_Q5.png}
	\caption{Q5 result: average of 127.42 routes per airline}
\end{figure}

\subsubsection*{Q6: What is the ISO code of the countries that have routes with an empty codeshare?}

\par
This query filters \texttt{routes} where \texttt{codeshare} is null or an empty
string, then joins the distinct source countries against the \texttt{countries}
table to obtain the ISO code.

\par
My first pass gave different counts: 208 rows from the DataFrame API, 207 from
SQL. Tracking it down, I found a real duplicate in \texttt{countries.txt}:
India shows up twice with the same \texttt{name} and \texttt{iso\_code}, just a
different \texttt{dafif\_code}. In the DataFrame API version I'd only called
\texttt{.distinct()} on \texttt{source\_country} before the join, so joining
against two India rows created two duplicate outputs that never got cleaned
up. The SQL version ran \texttt{DISTINCT} after the join instead, on the full
\texttt{(source\_country, iso\_code)} pair, which caught the duplicate. Adding
a second \texttt{.distinct()} after the join fixed it, both now return 207
rows:

\begin{figure}[H]
	\centering
	\includegraphics[width=0.9\textwidth]{4_1_Q6.png}
	\caption{Q6 result: 207 countries with routes containing an empty codeshare
		(DataFrame API and SQL, matching after fixing a duplicate-row discrepancy)}
\end{figure}
